%%%BLOCK id=036-03 ch=07 sec=碰撞·完全弹性碰撞 kind=formula alt=none cont=none y=0.41-0.55
\fbox{\textbf{完弹过程}}\quad 由动量守恒、能量守恒\quad 完弹$=$完非弹加反冲（有缓冲装置如弹簧、斜面）
\[
\left\{\begin{aligned}
m_1V_{01}+m_2V_{02}&=m_1V_1+m_2V_2\\
\tfrac12m_1V_{01}^2+\tfrac12m_2V_{02}^2&=\tfrac12m_1V_1^2+\tfrac12m_2V_2^2
\end{aligned}\right.
\qquad e=\frac{V_{01}-V_{02}}{V_2-V_1}=1 % CHECK: 恢复系数通常写作 (V2-V1)/(V01-V02)，原稿分子分母似为倒置(e=1时不影响)，照原样转录
\]
\[
\begin{aligned}V_1&=2V_{\text{共}}-V_{01}\\V_2&=2V_{\text{共}}-V_{02}\end{aligned}
\qquad\text{二倍共速减初速}
\]
\unclear{特}\ 动碰静
\[
V_1=\frac{m_1-m_2}{m_1+m_2}V_0\ \text{（动）}\qquad V_2=\frac{2m_1V_0}{m_1+m_2}\ \text{（静）}\qquad V_A+V_A'=V_B+V_B'
\]
%%%END
%%%BLOCK id=036-04 ch=07 sec=碰撞·完全非弹性碰撞 kind=formula alt=06,09,11,12 cont=none y=0.555-0.735
\fbox{\textbf{完非弹过程}}
\[
\left\{\begin{aligned}
m_1V_{01}+m_2V_{02}&=(m_1+m_2)V_{\text{共}}\qquad\text{共速}\\
\tfrac12(m_1+m_2)V_{\text{共}}-\Bigl(\tfrac12m_1V_{01}^2+\tfrac12m_2V_{02}^2\Bigr)&=\Delta E_K=E_{\text{某能}\max} % CHECK: 原稿 V共 处似无平方；E 下标“某能max”字迹难辨
\end{aligned}\right.
\]
%FIG p036_c 0.125 0.634 0.46 0.700
\notefig[0.45]{figures/p036_c.png}
%FIG p036_d 0.48 0.632 0.995 0.709
\notefig[0.7]{figures/p036_d.png}
\begin{align*}
\Delta E_K&=-\frac12\,\frac{m_1m_2}{m_1+m_2}(V_{01}-V_{02})^2=-E_{p\text{弹}\max}=-E_{p\unclear{G}\max}\\
&=Q=|fd|=Q_{\text{电}}=E_{\text{电场能}}\,.
\end{align*}
%%%END
%%%BLOCK id=036-05 ch=07 sec=碰撞·非完全弹性碰撞 kind=formula alt=none cont=none y=0.737-0.853
\fbox{\textbf{非完弹}}
\begin{align*}
m_1V_{01}+m_2V_{02}&=m_1V_{01}'+m_2V_{02}'\qquad\text{代入解方程}\\
\tfrac12m_1V_{01}^2+\tfrac12m_2V_{02}^2&=\tfrac12m_1V_{01}'^2+\tfrac12m_2V_{02}'^2+E_{\text{损}}
\end{align*}
%%%END
%%%BLOCK id=036-06 ch=07 sec=反冲·爆炸 kind=formula alt=none cont=none y=0.86-1.0
\fbox{\textbf{反冲}}
%FIG p036_e 0.07 0.878 0.36 0.929
\notefig[0.4]{figures/p036_e.png}
完弹$=$完非弹$+$反冲
\begin{align*}
m_1V_{01}-m_2V_{02}&=0 \\ % 第一项 V 的下标原稿似为“a”，按 V_{01} 转录
E&=\tfrac12m_1V_{01}^2+\tfrac12m_2V_{02}^2\\
V_{01}&=\sqrt{\frac{2m_2E}{m_1(m_1+m_2)}}\qquad V_{02}=\sqrt{\frac{2m_1E}{m_2(m_1+m_2)}}\\
E_{K01}&=\tfrac12m_1V_{01}^2=\frac{m_2}{m_1+m_2}E\qquad E_{K02}=\tfrac12m_2V_{02}^2=\frac{m_1}{m_1+m_2}E % 原稿 E_{K02} 式中 m 与下标2的写法为“½ ₂mV₀₂²”，按 m_2 转录
\end{align*}
%%%END
%%%BLOCK id=037-01 ch=07 sec=弹簧板块模型·压缩与反冲 kind=derivation alt=06 cont=none y=0.01-0.21
\textbf{反冲}\quad 势能$=$摩擦热
%FIG p037_a 0.10 0.058 0.265 0.108
\notefig[0.3]{figures/p037_a.png}

压缩、完非弹
\[
\left\{\begin{aligned}
m_1V_0&=(m_1+m_2)V_{\text{共}}\\
\tfrac12(m_1+m_2)V_{\text{共}}^2-\tfrac12m_1V_0^2&=-\bigl(E_{p\text{弹}}+Q_{\text{来}}\bigr)
\end{aligned}\right.
\]

反冲后又共速
\[
\left\{\begin{aligned}
(m_1+m_2)V_{\text{共}}&=(m_1+m_2)V_{\text{共}1}\\
\tfrac12(m_1+m_2)V_{\text{共}1}^2-\tfrac12(m_1+m_2)V_{\text{共}}^2&=E_{p\text{弹}}-Q_{\text{返}}=0\ \Rightarrow\ E_{p\text{弹}}=Q_{\text{返}}
\end{aligned}\right.
\]
\[
\frac12\,\frac{m_1m_2}{m_1+m_2}V_0^2=Q_{\text{来}}+Q_{\text{返}}=E_{p\text{弹}}+Q
\]
二次共速，反冲\quad 释放$=$消耗
%%%END
%%%BLOCK id=037-02 ch=07 sec=碰撞·拆分过程 kind=method alt=06 cont=none y=0.215-0.36
拆分过程。一个过程一组动量守恒、能量守恒方程
\begin{itemize}
\item 状态量：$P$\quad $E_K$\quad $F$
\item 过程量：$I$\quad $W$
\end{itemize}

完非弹$=$完非弹$+$反冲 % CHECK: 首项原稿似为“完非弹”，按字形照录(或应为“非完弹”)

完弹$=$完非弹$+$反冲

% 以下公式原稿被红笔框起
\begin{keybox}
\[
\Delta E_K=\frac12\,\frac{m_1m_2}{m_1+m_2}(V_{01}-V_{02})^2(e^2-1)
\]
\end{keybox}
（有\unclear{恢复}系数时）

每个过程只有2个体系，未参与该过程的物体原速不变
%%%END
%%%BLOCK id=037-03 ch=07 sec=多过程碰撞·双车与摆球 kind=problem alt=06,04 cont=none y=0.365-0.66
\begin{problem}[{[1]}]
%FIG p037_b 0.085 0.372 0.318 0.462
\notefig[0.4]{figures/p037_b.png}
\begin{solution}
\textbf{过程一}

①完非弹 $AB$
\[
mV-mV=2mV_{\text{共}},\qquad V_{\text{共}}=0 % CHECK: 第一个减号原稿字形近似“=”，按物理意义(两车动量反向)转为减号
\]
$AB$ 碰后速度为0
\[
\tfrac12\,2mV_{\text{共}}^2-\tfrac12mV^2-\tfrac12mV^2=\Delta E_{K1}
\]
\textbf{过程二}

②$C$ 与 $AB$ 完非弹
\begin{align*}
\frac m4V&=\Bigl(2m+\frac m4\Bigr)V_{\text{共}}'\\
\tfrac12\Bigl(2m+\frac m4\Bigr)V_{\text{共}}^2-\tfrac12\,\frac m4V^2&=\Delta E_{K2}=-\frac m4gh_{\max}\\
&\qquad\unclear{-\tfrac12\cdot2m} % 原稿第二式下方有叠写涂改的小字，形如“-½·2m”，位置在 ΔE_{K2} 下方，辨认不清
\\
&=\tfrac12\,\frac{\frac m4\cdot2m}{\frac m4+2m}(\Delta V_0)^2 % CHECK: 此行与上一行 ΔE_{K2} 的符号不一致，原稿行首为“=”，照录
\end{align*}
%FIG p037_c 0.28 0.59 0.44 0.66
\notefig[0.3]{figures/p037_c.png}
\[
\cos\theta=\frac{L-h_{\max}}{L}
\]
\end{solution}
\end{problem}
%%%END
%%%BLOCK id=037-04 ch=07 sec=多过程碰撞·ABC弹簧 kind=problem alt=06 cont=none y=0.67-1.0
\begin{problem}[{[2]}]
%FIG p037_d 0.10 0.70 0.32 0.74
\notefig[0.4]{figures/p037_d.png}
\begin{solution}
①$AB$ 爆炸反冲
\[
\left\{\begin{aligned}
m_AV_A-m_BV_B&=0 % 原稿 V_B 右上似有一小撇
\\
\tfrac12m_AV_A^2+\tfrac12m_BV_B^2&=E
\end{aligned}\right.
\qquad
V_A=\sqrt{\frac{2m_BE}{m_{\text{总}}m_A}},\quad V_B=\sqrt{\frac{2m_AE}{m_{\text{总}}m_B}}
\]
②$B$ 压 $C$ 完非弹
\[
\left\{\begin{aligned}
m_BV_B&=(m_B+m_C)V_{\text{共}1}\\
\tfrac12(m_B+m_C)V_{\text{共}1}^2-\tfrac12m_BV_B^2&=-E_{p\text{弹}\max}
\end{aligned}\right.
\]
\[
E_{p\text{弹}}=-\frac12\,\frac{m_Bm_C}{m_B+m_C}V_B^2 % CHECK: 原稿等号后有负号(弹性势能应为正值)，照录
\]
③$BC$ 完弹$=$完非弹$+$反冲
\[
\left\{\begin{aligned}
m_BV_B&=m_BV_B'+m_CV_C'\\
\tfrac12m_BV_B^2&=\tfrac12m_BV_B'^2+\tfrac12m_CV_C'^2
\end{aligned}\right.
\qquad
V_B'=\frac{m_B-m_C}{m_B+m_C}V_B,\quad V_C=\frac{2m_B}{m_B+m_C}V_B % CHECK: 此处 V_C 原稿无撇
\]
④$A$ 追上 $B$，$AB$ 完非弹
\[
\left\{\begin{aligned}
m_AV_A-m_BV_B'&=(m_A+m_B)V_{\text{共}}\\
\tfrac12(m_A+m_B)V_{\text{共}2}^2-\tfrac12m_AV_A^2-\tfrac12m_BV_B'^2&=\Delta E_{p\text{弹}}\,(-1)
\end{aligned}\right.
\]
\[
E_{p\text{弹}}=\frac12\,\frac{(m_A+m_B)m_C}{m_A+m_C+m_B}\bigl(V_{\text{共}2}-V_C'\bigr)^2
\]
\end{solution}
\end{problem}
%%%END
%%%BLOCK id=038-01 ch=07 sec=动量·不等式关系 kind=method alt=none cont=none y=0.0-0.085
% 原稿此块在页面右上角，被一条框线围起，框线下沿从第三行文字中间穿过
刚好 $\in$（上一次，下一次）

动量的不等式关系

推冰块8次完成，则 \unclear{$P\ge P_7$}，$P<P_9$ % 第一个 P 的下标被框线遮挡，符号“≥”处似箭头
%%%END
%%%BLOCK id=038-02 ch=07 sec=碰撞·碰撞发生条件 kind=knowledge alt=none cont=none y=0.05-0.205
\textbf{碰撞发生条件}\quad $V_{\text{后}}>V_{\text{前}}$ 撞得上\quad $V_{\text{后}}\le V_{\text{前}}$ 不穿越
\begin{itemize}
\item 动量守恒：$P_{01}+P_{02}=P_1+P_2$
\item 动能不增加：$\dfrac{P_{01}^2}{2m_1}+\dfrac{P_{02}^2}{2m_2}\ge\dfrac{P_1^2}{2m_1}+\dfrac{P_2^2}{2m_2}$ % 不等号处原稿有涂改叠写，辨为“≥”
\item 不穿越，不\unclear{击}穿：碰后 $V_{\text{后}}\le V_{\text{前}}$
\end{itemize}
\[
0\le e=\frac{V_2-V_1}{V_{01}-V_{02}}\le1
\]
%FIG p038_a 0.55 0.118 0.702 0.192
\notefig[0.3]{figures/p038_a.png}
\[
\left\{\begin{aligned}
V_{\text{完弹}}&\le V_1\le V_{\text{共}}\\
V_{\text{共}}&\le V_2\le V_{\text{完弹}}
\end{aligned}\right.
\]
%%%END
%%%BLOCK id=038-03 ch=07 sec=碰撞·传递系数 kind=formula alt=06 cont=none y=0.215-0.35
\textbf{速度传递}
\[
k_V=\frac{V_2}{V_0}=\frac{2m_1}{m_1+m_2}=\frac{2}{1+\dfrac{m_2}{m_1}}\ \to\ 2
\]
$m_2\ll m_1$ 时大球撞小球，小球传速\unclear{约}\red{为}\,\red{\unclear{走}} % 红笔：“为”字上有红色叉划，其后有红字，形如“走”，辨认不清

\textbf{动量传递}
\[
k_p=\frac{m_2V_2}{m_1V_0}=\frac{2m_2}{m_1+m_2}=\frac{2}{\dfrac{m_1}{m_2}+1}\ \to\ 2
\]
$m_1\ll m_2$ 时小球撞\red{大}墙，原速\red{$\times$}反\red{弹} % 红笔：“大”字上有红色叉划；“原速”后有红色斜线划去的一字；“反”后红笔补“弹”

\textbf{动能传递}
\[
k_{E_k}=\frac{\frac12m_2V_2^2}{\frac12m_1V_0^2}=k_p\cdot k_V=\frac{4m_1m_2}{(m_1+m_2)^2}=\frac{4}{\dfrac{m_1}{m_2}+\dfrac{m_2}{m_1}+2}
\]
$m_1=m_2$ 时等质换速
%%%END
%%%BLOCK id=038-04 ch=07 sec=动量守恒·适用条件 kind=knowledge alt=none cont=none y=0.36-0.455
动量守恒\quad $F_{\text{合}}=0$
\begin{itemize}
\item 内力极大动量守恒
\item 时间极短动量守恒
\item 水平/竖直方向动量守恒
\item 某阶段 $F$-$t$ 图动量守恒
\end{itemize}
$\Sigma I=mV_2-mV_1$ 可间接求冲量
%%%END
%%%BLOCK id=038-05 ch=07 sec=连续体冲量动量 kind=method alt=none cont=none y=0.38-0.62
\textbf{连续体冲量动量}\quad 流量 $Q=SV$.

\textbf{液柱类}
%FIG p038_b 0.125 0.434 0.305 0.50
\notefig[0.3]{figures/p038_b.png}
$V$ 为相对速度，
\[
-F\Delta t=\rho SV\Delta t\,(V_1-V_0)\qquad V_1=0\text{ 时 }F=-\rho SV^2.
\]
水平方向动量守恒
%FIG p038_c 0.09 0.512 0.256 0.612
\notefig[0.3]{figures/p038_c.png}
溅开
\[
F\Delta t=\rho SV_1\Delta t\,\bigl(-V_2\sin30^\circ-V_1\bigr)
\]
%%%END
%%%BLOCK id=038-06 ch=07 sec=连续体冲量动量 kind=method alt=04 cont=none y=0.60-0.78
\textbf{顶物并散开}\quad （此处不是水柱上升最高处，最高处速度为0，无法提供力）
%FIG p038_d 0.10 0.625 0.272 0.775
\notefig[0.3]{figures/p038_d.png}
\begin{align*}
V^2-V_0^2&=-2gh\\
F\Delta t&=\rho\,\Delta S\,\Delta l\,V\qquad F=\rho\,\Delta S\,V\,\frac{\Delta l}{\Delta t}\\
\Sigma\Delta F&=\rho SV\cdot V_0
\end{align*}
%%%END
%%%BLOCK id=038-07 ch=07 sec=连续体冲量动量 kind=method alt=06 cont=none y=0.78-0.93
单位体积有$n$个太空垃圾\,\red{（粘在飞船上）}\quad（完非弹） % 红笔波浪线划在“(粘在飞船上)”下
%FIG p038_e 0.06 0.795 0.22 0.86
\notefig[0.3]{figures/p038_e.png}
\[
FV\Delta t=\tfrac12\Delta m\,V^2
\]
\begin{keybox}
\[
F=\tfrac12\,mSV^2\cdot n
\]
\end{keybox}
% 上式原稿被红笔框起
\[
\bar F\,(V\Delta t)=\tfrac12\,(nmSV\Delta t)\,V^2
\]
%FIG p038_f 0.768 0.78 0.955 0.872
\notefig[0.35]{figures/p038_f.png}
或 $\bar F=\dfrac{\Delta\unclear{p}}{\Delta t}=\dfrac{m\cdot nSV\Delta t\cdot\frac V2}{\Delta t}=nmSV^2\cdot\dfrac12$\quad （碰上就停了，用$\bar V$算）。
%%%END
%%%BLOCK id=038-08 ch=07 sec=连续体冲量动量 kind=method alt=02 cont=none y=0.85-1.0
\textbf{双脚类}\quad ↓相对速度
%FIG p038_g 0.06 0.868 0.23 0.938
\notefig[0.3]{figures/p038_g.png}
\[
(m+M)g\Delta t=2\cdot\rho SV\Delta t\,V
\]
\[
\therefore\ V=\sqrt{\frac{(M+m)g}{2\rho S}}
\]
\[
\underset{\text{功率}}{P}=FV=(M+m)g\sqrt{\frac{(M+m)g}{2\rho S}}
\]
%%%END
%%%BLOCK id=039-01 ch=07 sec=连续体冲量动量 kind=problem alt=none cont=none y=0.07-0.21
每秒$N$粒米，每粒$m$，求$t$s后称读数 % 首字“每”被页眉“Date”字样压住；“称读数”三字笔迹有涂改叠写
\begin{gather*}
(F-N\Delta t\,mg)\Delta t=N\Delta t\,m\cdot V\qquad V=\sqrt{2gh}\\
N\Delta t\,mg\cdot\Delta t\approx0\qquad F=NmV\\
\text{故 }F_{\text{示数}}=F+G=Nm\sqrt{2gh}+N\unclear{\Delta t}\,mg % 末项中 Δt 处原稿为一团黑色涂迹
\end{gather*}
%%%END
%%%BLOCK id=039-02 ch=07 sec=连续体冲量动量 kind=problem alt=16 cont=none y=0.22-0.36
气体分子弹性撞器壁，$M$为单个重量，每立方米有$n_0$个，撞击时垂直分速度为$V$，求$P$压强 % “弹性”二字下有波浪线；M 原稿为大写
\[
\bar F=\frac{n_0mSV\Delta t\,V}{\Delta t}=\boxed{n_0mSV^2}\qquad F=\frac{m\Delta V}{\Delta t} % 原稿 n_0mSV^2 被圈出
\]
\[
P=\frac{\bar F}{S}=n_0mV^2
\]
%%%END
%%%BLOCK id=039-03 ch=07 sec=连续体冲量动量 kind=problem alt=06 cont=none y=0.365-0.77
\textbf{链条落地}
%FIG p039_a 0.055 0.389 0.16 0.47
\notefig[0.25]{figures/p039_a.png}
线密度 $\rho=\dfrac mL$

①落$x$时，求$N_{\text{地}}$
\begin{gather*}
(F-\rho\Delta x\,g)\Delta t=\Delta m\,V=\rho V\Delta x\\
\rho\Delta x\,g\Delta t\approx0\\
F=\rho V\frac{\Delta x}{\Delta t}=\rho V^2\qquad V^2=2gx\qquad\text{故 }F=2\rho gx\\
N=\underset{F}{2\rho gx}+\underset{G}{\rho gx}=\boxed{\frac{3mgx}{L}} % 原稿 3mgx/L 被圈出；2ρgx 下标 F，ρgx 下标 G
\end{gather*}
②求最大压力
\begin{gather*}
(F-\Delta m\,g)\Delta t=\rho V\cdot\Delta x\qquad\Delta m\,g\Delta t\approx0\\
\text{故 }F=\rho V\frac{\Delta x}{\Delta t}=\rho V^2\qquad V^2=2gL\\
F=2\rho gL\\
N=F+G=2\rho gL+\rho gL=3\rho gL=\boxed{3mg} % 原稿 3mg 被圈出
\end{gather*}
%%%END
%%%BLOCK id=063-04 ch=07 sec=动量·冲量 kind=formula alt=none cont=none y=0.43-0.53
\textbf{动量}
\begin{itemize}
\item 动量\quad 冲量
  \[
  \begin{array}{ll}
  p=mv\;=\;Ft & \qquad \unclear{…}\;F=\dfrac{\Delta p}{\Delta t}=ma \\[6pt]
  \phantom{p}=\sqrt{2mE_k} & \qquad E_k=\dfrac{p^2}{2m}
  \end{array}
  \]
\item 动量定理\quad $Ft=p_2-p_1=mv_2-mv_1$
\end{itemize}
%%%END
%%%BLOCK id=063-05 ch=07 sec=动量定理 kind=summary alt=none cont=none y=0.55-0.70
\begin{itemize}
\item 动量与冲量的理解\quad 矢量、合成分解. 以地为参 % CHECK: 末句可能为“以地为参考系”，原稿止于“参”
\item 动量定理的理解与应用
  \begin{itemize}
  \item $I=Ft$
  \item 图像法\quad $F$-$t$图
  \item 平均方法\quad $I=\dfrac{F_1+F_2}{2}\,t$\quad 随时间 $F=kt$
  \item 动量定理法\quad $I=Ft=\Delta p$ （自 $\Delta p$ 处画一条弧形箭头指回 $I$）
  \end{itemize}
\end{itemize}
%%%END
%%%BLOCK id=063-06 ch=07 sec=动量守恒 kind=summary alt=none cont=none y=0.69-0.80
\textbf{动量守恒定律及应用}
\begin{itemize}
\item 动量守恒定律\quad (合力为0)
\item 碰撞、爆炸、反冲.
\end{itemize}
%%%END
%%%BLOCK id=063-07 ch=07 sec=碰撞 kind=summary alt=none cont=none y=0.81-0.95
\textbf{动量守恒定律的理解与应用}
\begin{itemize}
\item 碰撞问题.
  \begin{itemize}
  \item 动量守恒，动能不增加，速度合理，不穿越
  \end{itemize}
\end{itemize}
%%%END
%%%BLOCK id=063-08 ch=07 sec=爆炸反冲·人船模型 kind=summary alt=none cont=none y=0.95-1.00
% 页面底边被照片裁切，下方内容不可见
\begin{itemize}
\item 爆炸和反冲问题
  \begin{itemize}
  \item 人船模型\quad 人走船走\ 人停船停
  \end{itemize}
\end{itemize}
%FIG p063_a 0.05 0.968 0.235 1.0
\notefig[0.25]{figures/p063_a.png}
% 图为页底左侧几幅被裁切的小示意图（一横线、带斜面的块、弧形槽、带球小车）
\[ x_{\text{人}}=\dfrac{m_{\text{船}}}{m_{\text{人}}+m_{\text{船}}}L_{\text{相}} \]
\[ x_{\text{船}}=\dfrac{m_{\text{人}}}{m_{\text{人}}+m_{\text{船}}}L_{\text{相}} \] % CHECK: 页底裁切，第二式分母及 L 的下标（可能为“相对”）据残笔读出
%%%END
%%%BLOCK id=064-01 ch=07 sec=子弹打木块模型 kind=derivation alt=none cont=none y=0.02-0.18
\textbf{动量守恒的四类模型}

\textbf{子弹打木块 模型}
%FIG p064_a 0.313 0.027 0.585 0.098
\notefig[0.4]{figures/p064_a.png}
% 图中文字：水平面光滑、m、v_0（箭头）、M、d、f
\textbf{未穿出}
\begin{align*}
mv_0&=(M+m)V, & V&=\frac{mv_0}{M+m}\\
ft&=MV-0, & t&=\frac{Mmv_0}{f(M+m)}\\ % 原稿 ft 之前有一淡痕（疑为背面透字）
-fx_1&=\tfrac12 mV^2-\tfrac12 mv_0^2, & x_1&=\frac{Mm(M+2m)v_0^2}{2f(M+m)^2}\\ % CHECK: 分子括号内首字母 M/m 难分，按推导读作 M+2m
fx_2&=\tfrac12 MV^2, & x_2&=\frac{Mm^2v_0^2}{2f(M+m)^2}
\end{align*}
\[ \text{进入深度}\quad x_{\text{相对}}=x_1-x_2=\frac{Mmv_0^2}{2f(M+m)} \]
%%%END
%%%BLOCK id=064-02 ch=07 sec=子弹打木块模型 kind=formula alt=none cont=none y=0.18-0.28
% 原稿：自上方“未穿出”推导区右下引出一条弧线并带箭头向下，指向右侧“若子弹穿出速度为 v_1”部分
\textbf{未穿出}
\begin{align*}
mv_0&=(m+M)V\\
Q_{\text{热}}&=fx_{\text{相对}}=\tfrac12 mv_0^2-\tfrac12(M+m)V^2
\end{align*}
\textbf{穿出}
\begin{align*}
mv_0&=mv_1+MV_2\\
Q_{\text{热}}&=fL=\tfrac12 mv_0^2-\tfrac12 mv_1^2-\tfrac12 MV_2^2
\end{align*}
%%%END
%%%BLOCK id=064-03 ch=07 sec=子弹打木块模型 kind=derivation alt=none cont=none y=0.28-0.53
\textbf{若子弹穿出速度为 $v_1$}
\begin{align*}
mv_0&=mv_1+MV_2\\
V_2&=\frac{m(v_0-v_1)}{M}
\end{align*}
\begin{itemize}
\item 对子弹\quad $W_1=-fx_1=\tfrac12 m(v_1^2-v_0^2)$
\item $W_2=$ 对块\quad $fx_2=\tfrac12 MV_2^2$
\end{itemize}
\[ \therefore\ W_2=\frac{m^2(v_0-v_1)^2}{2M} \]
\[ Q=fx_1-fx_2=fd. \]
%%%END
%%%BLOCK id=064-04 ch=07 sec=弹簧模型 kind=method alt=none cont=none y=0.27-0.41
\textbf{物体弹簧 模型\quad 双振子}
%FIG p064_b 0.14 0.306 0.37 0.395
\notefig[0.3]{figures/p064_b.png}
% 图中文字：v、t（v-t 图像，两条曲线与若干虚线、竖线）
\begin{itemize}
\item 有完弹\quad 有完非弹.
\item \unclear{共}速为压缩最大\quad 完非弹 % CHECK: 首字形似“失”，按语义读作“共”
\item 伸开时伸长最大\quad 完弹. % CHECK: “伸开时伸长”字迹潦草，读法存疑
\end{itemize}
%%%END
%%%BLOCK id=064-05 ch=07 sec=滑块曲面体模型 kind=method alt=none cont=none y=0.41-0.54
\textbf{滑块曲面体 模型}\quad 水平方向动量守恒
\[ \left\{\begin{array}{l}\text{上去是完非弹}\\[2pt] \text{下去是完弹}\end{array}\right. \]
%%%END
%%%BLOCK id=064-06 ch=07 sec=滑块滑板模型 kind=method alt=none cont=none y=0.54-0.66
\textbf{滑块滑板模型}\quad 地光滑时动量守恒
\begin{itemize}
\item 板长够共速为完非弹
\item 板长不够共速为子弹打木块.
\end{itemize}
%%%END
%%%BLOCK id=067-01 ch=07 sec=力学三大观点综合应用 kind=method alt=03,04,06 cont=none y=0.07-0.50
\textbf{力学三大观点的综合应用}

直线运动、平抛运动、圆周运动的组合。

动力学、能量、动量的综合应用

要列各物理量在某一时刻的关系式，用牛二

研究某一物体受到力的持续作用发生运动状态改变\quad 用动量定理\ （涉及时间）\ 或动能定理\ （涉及位移） % 原稿“(涉及时间)”“(涉及位移)”写在本行上方，分别位于“用动量定理”“或动能定理”之上

研究系统，用动量守恒定律和机械能守恒定律\quad 注意守恒的条件

涉及相对位移，优先考虑能量守恒定律，系统克服摩擦力所做的总功等于系统机械能的减少量（转变为系统内能的量）。

在涉及碰撞、爆炸、打击、绳绷紧等现象时，一般隐含有系统机械能损失，与其他形式的能量的转换，作用时间都极短，用动量守恒定律解决。
%%%END
%%%BLOCK id=074-01 ch=07 sec=子弹打木块·多块穿透 kind=problem alt=06 cont=none y=0.04-0.50
\begin{problem}[1、]
%FIG p074_a 0.075 0.05 0.285 0.135
\notefig[0.3]{figures/p074_a.png}
恰好射穿.

%FIG p074_b 0.075 0.135 0.31 0.205
\notefig[0.3]{figures/p074_b.png}
求进入第二块钢块的深度
\begin{solution}
\textbf{Step1.}\quad $mV_0=(2m+m)V$\quad $V=\dfrac{V_0}{3}$
\[
\Delta E=\frac{1}{2}mV_0^{2}-\frac{1}{2}\times 3mV^{2}=\frac{1}{3}mV_0^{2}
\]
\textbf{Step2.}\quad $mV_0=mV_1+mV_2$\quad 由 $f$ 为恒力 故射出一块\unclear{耗}能为 $\dfrac{\Delta E}{2}$
\begin{align*}
&\frac{1}{2}mV_0^{2}=\frac{1}{2}mV_2^{2}+\frac{1}{2}mV_1^{2}+\frac{1}{2}\Delta E\qquad \text{因 }V_1>V_2\quad \text{得 }V_1=\Bigl(\frac{1}{2}+\frac{\sqrt{3}}{6}\Bigr)V_0\\
&mV_1=2mV'\\
&\Delta E'=\frac{1}{2}mV_1^{2}-\frac{1}{2}\,2mV'^{2}=\frac{1}{2}\Bigl(1+\frac{\sqrt{3}}{2}\Bigr)\times\frac{1}{2}\Delta E\\
&\Delta E'=fx\\
&\frac{1}{2}\Delta E=fd\qquad \therefore\ x=\frac{1}{2}\Bigl(1+\frac{\sqrt{3}}{2}\Bigr)d
\end{align*}
\end{solution}
\end{problem}
%%%END
%%%BLOCK id=074-02 ch=07 sec=动量定理·柱状模型 kind=method alt=none cont=none y=0.54-0.78
\textbf{两类柱状模型}
\begin{itemize}
\item \textbf{连续性}\quad 流体类、液流气流.\quad $\Delta m=\rho Sv\Delta t$\quad （此式上方批注：给密度$\rho$；页面右侧批注：注意 相对速度 与速度分解；分方向动量定理.）
  \begin{itemize}
  \item 用动量定理研究这段流体
  \item $F\Delta t=\rho Sv\Delta t(V_{\text{末}}-V_0)$\qquad $F=\rho Sv(V_{\text{末}}-V_0)$\quad 注意正负
  \end{itemize}
\item \textbf{粒子性}\quad 微粒类、电子流、光子流、尘埃、给出单体积内粒子数
  \begin{itemize}
  \item 微元内粒数 $N=nv_0S\Delta t$,
  \item 用动量定理研究单个粒子，建立方程 再乘$N$计算
  \end{itemize}
\end{itemize}
%%%END
%%%BLOCK id=074-03 ch=07 sec=动量定理·柱状模型 kind=derivation alt=09,17 cont=none y=0.83-1.00
\textbf{带压强}\quad $F\Delta t=\Delta mV=\rho\Delta S\,hV$
\[
P_{\text{压}}=\frac{F}{\Delta S}=\frac{\rho hv}{\Delta t}
\]
\[
V_{\text{增}}=\frac{h}{\Delta t}=\frac{P}{\rho V} % CHECK: 下标字迹涂改，疑为「增」
\]

\textbf{带电流}
\[
qU=\frac{1}{2}mv^{2}
\]
\[
Q=I\Delta t
\]
\unclear{总}\ $M=\dfrac{Q}{q}m=\dfrac{I\Delta t}{q}m$
\[
\bar F\Delta t=Mv
\]
\[
\bar F=I\sqrt{\frac{2mU}{q}}
\]

\textbf{带光电子} % CHECK: 标题字迹似「带光电子」，但下方内容全是光子（ε=hν、p=h/λ），疑为「带光子」，按原样转录
\[
\varepsilon=h\nu\qquad p=\frac{h}{\lambda}\qquad c=\lambda\nu
\]
\[
\therefore\ \varepsilon=pc
\]
\[
E=\frac{n}{t}\,\frac{\varepsilon}{S}\qquad \unit{J/s/m^{2}}
\]
设$t$时间$n$个光子 每个\unclear{粒子}（下一行）为$\varepsilon$ % 此句写在上式右侧，「为ε」另起一行

\unclear{全部反射}\quad $Ft=2p\times n$

$\Delta p=2p_0$\quad $F=ma$
%%%END
%%%BLOCK id=079-01 ch=07 sec=动量定理·系统动量定理 kind=knowledge alt=04 cont=none y=0.04-0.24
\textbf{系统动量定理}\quad $\sum\vec I=\sum\vec F\,t=m\vec V_2-m\vec V_1$

%FIG p079_a 0.0 0.13 0.19 0.22
\notefig[0.25]{figures/p079_a.png}
\[
mgt=m\Delta V
\]
%FIG p079_b 0.36 0.13 0.505 0.22
\notefig[0.25]{figures/p079_b.png}
\begin{align*}
&mg\,2R=\tfrac12 mV^2-\tfrac12 mV_0^2\\
&I_N+mgt=m\vec V-m\vec V_0\\
&W_N=0\qquad I_N\ne0
\end{align*}
%%%END
%%%BLOCK id=079-02 ch=07 sec=动量定理·系统动量定理 kind=problem alt=03 cont=none y=0.26-0.35
\begin{problem}
$M$带$m$、$V_0$时$M$离开 $m$停后 求$V_M$
%FIG p079_c 0.02 0.30 0.17 0.335
\notefig[0.25]{figures/p079_c.png}
\begin{solution}
\[
\begin{cases}
(M+m)at=MV+m\cdot0-(M+m)v\\[2pt]
t=\dfrac{V_0}{\mu g}
\end{cases}
\]
% CHECK: t 的分母手写像 ug，按 \mu g 转录
\end{solution}
\end{problem}
%%%END
%%%BLOCK id=079-03 ch=07 sec=动量守恒·适用条件 kind=knowledge alt=none cont=none y=0.33-0.42
% 原稿位置：页面右侧，与“M带m”“绳绷紧”两例同高；右侧另有一括号括出“分阶段”“分方向”
\[
\text{动量守恒}
\left\{
\begin{array}{l}
\text{内力}\gg\text{外力}\\
\text{时间极短}\\
\text{合外力为0}
\end{array}
\right.
\left\{
\begin{array}{l}
\text{分阶段}\\
\text{分方向}
\end{array}
\right.
\]
%%%END
%%%BLOCK id=079-04 ch=07 sec=动量定理·系统动量定理 kind=problem alt=none cont=none y=0.36-0.53
\begin{problem}
% 题干只有配图中的标注（$V_0=2\,\mathrm{m/s}$、$L=0.4\,\mathrm{m}$）
%FIG p079_d 0.005 0.365 0.195 0.495
\notefig[0.25]{figures/p079_d.png}
\begin{solution}
\[
-(M+m)g\,\Delta t=(m+M)V_{\text{共}}-(MV_0-mV_0)
\]
\[
\Delta t=\dfrac{0.4}{4}=0.1\,\mathrm{s}\quad {=}\,\dfrac{L}{V_{\text{相}}}
\]
% 原稿“=L/V相”写在“0.1s”右上方；CHECK: 0.1s 的 s 手写像 5
\[
\therefore\ V_{\text{共}}=\dfrac{M-3m}{M+m}\quad
\begin{cases}
M>3m\text{ 时 向上}\\
M=3m\text{ 时 不动}\\
M<3m\text{ 时 向下.}
\end{cases}
\]
\end{solution}
\end{problem}
%%%END
%%%BLOCK id=079-05 ch=07 sec=反冲·爆炸 kind=note alt=none cont=none y=0.44-0.51
喷气切记质量减少
% 句末有一斜划
\[
MV_0=(M-m)V_2+MV_1
\]
% CHECK: 末项质量手写像 M（按物理应为 m）
%%%END
%%%BLOCK id=079-06 ch=07 sec=滑块曲面体模型 kind=method alt=06 cont=none y=0.575-0.635
%FIG p079_e 0.02 0.575 0.125 0.632
\notefig[0.25]{figures/p079_e.png}
机守可滑到最高处\quad \underline{水平方向}动量守恒\quad 动量不守恒
%%%END
%%%BLOCK id=079-07 ch=07 sec=碰撞·非完全弹性碰撞 kind=problem alt=06 cont=none y=0.675-0.91
\begin{problem}
%FIG p079_f 0.04 0.675 0.215 0.745
\notefig[0.3]{figures/p079_f.png}
在 $L$ 区内 $A$ 受 $F$ 向右由静止匀加、与 $B$ 只碰 1 次 最终两\unclear{距}球距离不变为 $\dfrac49L$\\
求 $\dfrac{m_A}{m_B}$ 和 $E$机\unclear{损}
% “两”与“球”之间多一个像“距”的笔画；“E机损”末字像 损
\begin{solution}
\[
\begin{cases}
m_AV_0=m_BV_B-m_AV_A\\
V_A=V_B=V\\
FL=\tfrac12 m_AV_0^2\\
Ft=2m_AV\\
\dfrac49L=Vt
\end{cases}
\qquad\therefore\ \dfrac{m_A}{m_B}=\dfrac14
\]
\[
\Delta E_{\text{机}}=\tfrac12 m_AV_0^2-\left(\tfrac12 m_AV^2+\tfrac12 m_BV^2\right)=\dfrac49FL
\]
% CHECK: 括号内两项之间的符号手写像一个潦草的 +，按 + 转录
\end{solution}
\end{problem}
%%%END
%%%BLOCK id=079-08 ch=07 sec=滑块曲面体模型 kind=problem alt=06 cont=none y=0.91-1.00
\begin{problem}
%FIG p079_g 0.03 0.918 0.205 0.99
\notefig[0.25]{figures/p079_g.png}
C 在最低点时 $AB$ 分离\quad C 到 $A$ 左侧最高时 $V_A=\dfrac V4$ 向左．$A$、$B$ 会滑下
% CHECK: 若取向右为正，V_AC=V/4>0 应向右，手写为“向左”，原样转录
\begin{solution}
① $AB$ 共同向右\quad $mV=2mV_{AB}$\\
$mV+m\left(-\dfrac V2\right)=2mV_{AC}$\qquad $V_{AC}=\dfrac V4$\\
② $AC$ \unclear{左}减速，$AB$ 分开，$B$ 向右匀速
% 页面底部此行被截断，字迹只见上半部
\end{solution}
\end{problem}
%%%END
%%%BLOCK id=080-01 ch=07 sec=质心运动定理 kind=knowledge alt=03 cont=none y=0.00-0.06
% 页顶边距处一行；首行右端被页边裁切，“使质心静止或匀速直线”处字迹残缺，“运动”折到第二行
\textbf{质心运动定理}\quad 系统受力\quad 质心加速运动\quad $F_{\text{合}}=0$ 使\unclear{质}心\unclear{静止}或匀速直\unclear{线}运动
%%%END
%%%BLOCK id=080-02 ch=07 sec=弹簧模型·弹簧双振子 kind=method alt=06 cont=none y=0.06-0.20
\textbf{弹簧双振子模型}

%FIG p080_a 0.10 0.105 0.265 0.19
\notefig[0.28]{figures/p080_a.png}
%FIG p080_b 0.295 0.08 0.48 0.20
\notefig[0.3]{figures/p080_b.png}
$a=0$ 时 $V_{\max}$\\
质心向右匀速直线
% 句末有一小撇“'”
%%%END
%%%BLOCK id=080-03 ch=07 sec=弹簧模型·弹簧双振子 kind=method alt=06 cont=none y=0.18-0.29
%FIG p080_c 0.09 0.22 0.31 0.27
\notefig[0.3]{figures/p080_c.png}
%FIG p080_d 0.30 0.20 0.49 0.28
\notefig[0.28]{figures/p080_d.png}
一次周期中\\
$V_A=V_B$ 时\quad $V_{A\max}$\quad $V_{B\max}$\quad $E_{pT\max}$\quad $E_{\text{机}\max}$\\
$a_1=a_2$ 时\quad $\Delta V_{\max}$
% 原稿“V_A=V_B时”一行四项之间只有空格、无标点；E_p 的下标像 T
%%%END
%%%BLOCK id=080-04 ch=07 sec=弹簧模型·弹簧双振子 kind=problem alt=06 cont=none y=0.295-0.50
\begin{problem}
%FIG p080_e 0.09 0.295 0.32 0.365
\notefig[0.3]{figures/p080_e.png}
$W_F=W$ 突然撤 $F$．求 $I_{\text{墙}A}$、$V_A$ $V_{B\min}$
% CHECK: V_A 后似漏写 max（下文有 V_{A max}），原样转录
\begin{solution}
\begin{align*}
&W=E_p=\tfrac12 m_BV_0^2\qquad\therefore\ V_0=\sqrt{\dfrac{2W}{3m}}\qquad I_{\text{墙}A}=m_BV_0=\sqrt{6mW}\\
&m_BV_0=m_AV_{A\max}+m_BV_{B\min}\\
&\tfrac12 m_BV_0^2=\tfrac12 m_AV_{A\max}^2+\tfrac12 m_BV_{B\min}^2\\
&V_{B\min}=\dfrac{m_B-m_A}{m_B+m_A}V_0=\sqrt{\dfrac{W}{6m}}\qquad V_{A\min}=0
\end{align*}
\end{solution}
\end{problem}
%%%END
%%%BLOCK id=080-05 ch=07 sec=弹簧模型·弹簧双振子 kind=method alt=06 cont=none y=0.50-0.65
$F_1=F_2$\quad $W_{\text{机}}\uparrow$\quad $T=F_1=F_2$ 时 $E_k$ 最大\\
系统 $M+m$ 动量不为0
%FIG p080_f 0.04 0.51 0.31 0.63
\notefig[0.3]{figures/p080_f.png}
%FIG p080_g 0.31 0.558 0.51 0.65
\notefig[0.28]{figures/p080_g.png}
若 $m<M$
% 图中标注：S_1、F_1、F_2、S_2、μ=0、W_1>0、W_2>0；v-t 图纵轴标 m 与 -M
%%%END
%%%BLOCK id=080-06 ch=07 sec=弹簧模型·弹簧双振子 kind=summary alt=none cont=none y=0.69-1.00
\textbf{结论}

%FIG p080_h 0.03 0.708 0.68 0.83
\notefig[0.65]{figures/p080_h.png}
%FIG p080_i 0.025 0.824 0.72 0.988
\notefig[0.7]{figures/p080_i.png}
% 图内标注：上方左右两个装置图（左：m_1 带 V_0 向右；右：m_2 带 V_0 向右），各行左图对应左装置、右图对应右装置；行标 ① m_1=m_2（左图下方另标“原压原拉”）、② m_1>m_2、③ m_1<m_2；各曲线分别标 m_1、m_2
%%%END
%%%BLOCK id=089-08 ch=07 sec=人船模型 kind=method alt=none cont=none y=0.85-0.98
\noindent 人船模型
\[ m_1V_1=m_2V_2 \]
\[ m_1S_1=m_2S_2 \]
\[ S_1+S_2=S_{\text{相}} \]
\[ \begin{cases} S_1=\dfrac{m_2}{m_1+m_2}S_{\text{相}}\\[6pt] S_2=\dfrac{m_1}{m_1+m_2}S_{\text{相}}\end{cases} \]
\noindent 由牛一 $\Sigma F=0\ \longrightarrow$ 静止 或 匀速直线运动 \\
系统 $\Sigma F=0\ \longrightarrow$ 质心 静止 或 匀速直线运动
%%%END
%%%BLOCK id=091-11 ch=07 sec=碰撞·速度范围 kind=knowledge alt=none cont=none y=0.95-0.99
碰撞定义域\quad $v_{\text{完非}}\le v\le v_{\text{完弹}}$
%%%END
%%%BLOCK id=119-04 ch=07 sec=动量与能量综合·多球系统 kind=problem alt=06 cont=none y=0.63-0.79
\begin{problem}[2.]
%FIG p119_e 0.085 0.642 0.265 0.722
\notefig[0.25]{figures/p119_e.png}
% 图中标注：$m_1$、$M$、$m_2$（水平连线），$M$ 处向上箭头 $v_0$

求\ 1\ 2\ 相遇时速度
\begin{solution}
\begin{align*}
Mv_0&=(M+2m)v_y\\
\frac{1}{2}Mv_0^2&=\frac{1}{2}Mv_y^2+\frac{1}{2}\,2m\left(v_x^2+v_y^2\right)\\
v_x&=\sqrt{\frac{M}{M+2m}}\,v_0\qquad v_y=\frac{Mv_0}{M+2m}\\
v&=\sqrt{v_x^2+v_y^2}=\sqrt{\frac{2M(M+m)}{M+2m}}\,v_0 % CHECK: 根号是否只盖分子；若按 v_x、v_y 代入，分母应为 (M+2m) 而非根号下的 M+2m
\end{align*}
\end{solution}
\end{problem}
%%%END
%%%BLOCK id=120-01 ch=07 sec=动量与能量综合·多球系统 kind=problem alt=06,03 cont=none y=0.03-0.31
% 页面右上角露出邻页(化学)内容“…和CO2混合…×22.4=6.72L”，不属于本页，已忽略；页面下方为空白横线(背面透出淡淡反字，已忽略)
\begin{problem}
%FIG p120_a 0.025 0.035 0.148 0.128
\notefig[0.25]{figures/p120_a.png}
% 图中标注：A、B、C 三球在同一水平线上，B 处向上箭头 $v_0$，间距 $L$、$L$

① AC第1次碰\quad 求$v_B$

② 再次三球共线\ $v_B$、绳中拉力$F$

③ 运动过程中A最大动能和两绳夹角
\begin{solution}
① $mv_0=3mv_{\text{共}}$\quad $v_{\text{共}}=\dfrac{v_0}{3}$

② \[
\left\{\begin{aligned}
&mv_0=mv_B+2mv_A\\
&\frac{1}{2}mv_0^2=\frac{1}{2}mv_B^2+\frac{1}{2}m\left(v_A^2+v_A^2\right)
\end{aligned}\right.
\]
\[
\left\{\begin{aligned}
&v_B=-\frac{1}{3}v_0\\
&v_A=\frac{2}{3}v_0
\end{aligned}\right.
\qquad
F=\frac{m v_{\text{相}}^2}{L}=\frac{mv_0^2}{L}
\]
AC相对B\quad $v_{\text{相}}=v_0$

③
%FIG p120_b 0.65 0.118 0.83 0.17
\notefig[0.25]{figures/p120_b.png}
% 图中标注：$v_A$（左上箭头）、$v_C$（右上箭头）、45°，两处直角记号，顶点处有夹角记号

沿绳$v_B$为0，故AC均无分量\quad $E_{kA\max}$\quad $v_B=0$

$v_A$\ $v_C$垂绳\quad $\theta=90^\circ$
\[ 2m_Av\sin 45^\circ=mv_0 \]% CHECK: 下标 A 的位置疑应在 v 上(2mv_A sin45°=mv_0)
\[ v_A=v_C=\frac{\sqrt{2}}{2}v_0 \]
\[ E_{kA}=E_{kC}=\frac{mv_0^2}{4} \]
$\theta=90^\circ$ % 此 θ=90° 在页面中部偏右、字较大，疑为 ③ 的最终结论
\end{solution}
\end{problem}
%%%END
%%%BLOCK id=196-01 ch=07 sec=动量定理·板块变力冲量 kind=problem alt=03,06 cont=none y=0.00-0.30
% 原稿此块被一条自页面上方向左下弯曲、末端回钩的大弧线划过，疑似划去，仍照录
\begin{problem}
\notefig[0.28]{figures/p196_a.png}
\begin{solution}
\begin{align*}
&F=(M+m)a\\
&f=\mu Mg=Ma\qquad a=2\unit{m/s^2}\qquad F=6\unit{N}\qquad \struck{t=6\unit{s}}\qquad t=1\unit{s}\\
&\mu Mg=\frac{M}{M+m}F=\frac{2}{3}F=4 \quad % 分子 $M$ 左侧有一团涂改墨迹
\end{align*}
\[
F=6\unit{N}\qquad t=6\unit{s}\qquad \bar F=\frac{6}{2}\unit{N}=3\unit{N}
\]
% CHECK: 此处 $t=6\unit{s}$ 与后文 $t=1\unit{s}$（$\bar F t=3\unit{N\cdot s}$）不一致
\begin{align*}
&\underline{v_{\text{甲}}=v_{\text{乙}}=}\ \ \unclear{…}\ I_F=\bar F t=3\unit{N\cdot s}\\
&\bar F t=(M+m)v\qquad v=1\unit{m/s}\qquad \Delta P_{\text{乙}}=mv=1\unit{kg\cdot m/s}\\
&I_f=\Delta P_{\text{甲}}=Mv=2\unit{N\cdot s}\\
&W_{f\text{对甲}}=E_{k\text{甲}}=\tfrac{1}{2}mV^2=\unclear{1\unit{J}} % CHECK: 动能应为 $\tfrac12Mv^2$（$M=2\unit{kg}$）；末项手写形似“U”，按 $1\unit{J}$ 记
\end{align*}
\end{solution}
\end{problem}
%%%END
